\documentclass[12pt,a4paper]{article}
\usepackage{amsmath,amssymb,amsthm}
\usepackage{hyperref}
\usepackage{geometry}
\usepackage{enumitem}
\geometry{margin=2.5cm}

\newtheorem{theorem}{Theorem}[section]
\newtheorem{corollary}[theorem]{Corollary}
\theoremstyle{definition}
\newtheorem{definition}[theorem]{Definition}
\theoremstyle{remark}
\newtheorem{remark}[theorem]{Remark}

\title{Charge Quantization Without Monopoles:\\
Dimension-Forced Discrete Charges in Recognition Science}

\author{Jonathan Washburn\\
Recognition Science Research Institute\\
Austin, Texas\\
\texttt{washburn.jonathan@gmail.com}}

\date{March 2026}

\begin{document}
\maketitle

\begin{abstract}
Charge quantization in Recognition Science follows from the discrete
$D$-cube structure, not magnetic monopoles.  Face-pair count
$\mathrm{face\_pairs}(D) = D(D-1)/2 = 3$ for $D = 3$; charges are
integer-valued from the discrete ledger.  Dirac condition $qg \in
\mathbb{Z}$ holds trivially.  RS predicts no fundamental or GUT
monopoles; the gauge group is dimension-forced.  All searches (MoEDAL,
IceCube, MACRO) return null results.
\end{abstract}

\section{Introduction}
\label{sec:intro}

Dirac showed in 1931~\cite{Dirac1931} that a single magnetic monopole
would explain charge quantization: the electron wave function in a
monopole field must be single-valued, so the phase around a loop is
$2\pi n$, yielding $qg = n\hbar c/2$ ($n \in \mathbb{Z}$).  If $g \neq
0$, $q$ is quantized.  't Hooft~\cite{tHooft1974} and
Polyakov~\cite{Polyakov1974} showed that GUT symmetry breaking
$G_{\mathrm{GUT}} \to SU(3) \times SU(2) \times U(1)$ produces
't Hooft-Polyakov monopoles (mass $\sim 10^{15}$--$10^{16}$ GeV),
copiously produced in the early universe.  Despite decades of
searching, none have been found.  RS offers a different resolution:
charge quantization follows from the discrete ledger structure, not
monopoles.  RS predicts no fundamental or GUT monopoles exist.

\section{Charge Quantization from the Discrete Ledger}
\label{sec:charge-quantization}

In Recognition Science~\cite{RS_Axioms} (RS), physical states are
recorded on a double-entry recognition ledger---discrete,
lattice-indexed, with cost $J(x) = \tfrac{1}{2}(x + x^{-1}) - 1$.
Charge quantization follows.

\begin{definition}[Recognition Ledger]
Discrete structure recording physical states; allowed values form a
discrete set.
\end{definition}

\begin{theorem}[Charge Types from Face Pairs]
\label{thm:charge-types}
For spatial dimension $D$, the number of independent charge types is
\begin{equation}
\mathrm{charge\_types}(D) = \binom{D}{2} = \frac{D(D-1)}{2}.
\end{equation}
For $D = 3$, there are exactly three charge types (color, weak
isospin, hypercharge).
\end{theorem}

\begin{proof}
The three 2D coordinate planes of the $D$-cube (pairs of distinct
basis axes: $\{xy, xz, yz\}$) number $\binom{D}{2} = D(D-1)/2$.
For $D = 3$: planes $(xy)$, $(xz)$, $(yz)$ map to $SU(3)_C$,
$SU(2)_L$, $U(1)_Y$ respectively (three gauge sectors).
\end{proof}

\begin{remark}[Two Counting Conventions for ``Face Pairs'']
\label{rem:conventions}
In a $D=3$ cube, ``face pairs'' has two natural meanings: (i) pairs
of opposite faces, $F/2 = D = 3$; (ii) pairs of coordinate axes,
$\binom{D}{2} = 3$.  Both give 3 for $D = 3$; the $\binom{D}{2}$
convention is used here because it directly indexes the 2D gauge
planes $SU(3)_C$, $SU(2)_L$, $U(1)_Y$.
\end{remark}

\begin{theorem}[Discrete Charges from Ledger Structure]
\label{thm:discrete-charges}
Charges are quantized (integer-valued) because the ledger has discrete
entries.  Quantization does not require monopoles.
\end{theorem}

\begin{proof}
Ledger entries are countable; charge sectors are ledger channels with
integer indices (color $1,2,3$; weak isospin $\pm 1/2$; hypercharge).
Double-entry structure enforces conservation as sums of integer
contributions.  Discreteness suffices; no monopole invoked.
\end{proof}

\begin{corollary}[Dirac Condition Trivially Satisfied]
\label{cor:dirac}
$qg \in \mathbb{Z}$ holds trivially: all $q$ are integers.  The
condition imposes no constraint; quantization is ledger-enforced.
\end{corollary}

\begin{theorem}[Four Charge Types Excluded]
\label{thm:four-excluded}
$\mathrm{charge\_types} \neq 4$: $D(D-1)/2 = 4$ has no integer $D > 0$.
\end{theorem}

\begin{proof}
$D^2 - D - 8 = 0$ gives $D = (1 \pm \sqrt{33})/2 \notin \mathbb{Z}$.
\end{proof}

\begin{remark}
RS derives three gauge factors from $D = 3$; no free parameter.
\end{remark}

\section{Why GUT Monopoles Do Not Exist in RS}
\label{sec:no-gut-monopoles}

't Hooft-Polyakov monopoles arise when $G \to H$ with $\pi_2(G/H)
\neq 0$.  In $SU(5)$ GUT, breaking to the SM yields monopoles; same for
$SO(10)$, $E_6$.  RS does not employ a GUT~\cite{RS_Gauge}.  The gauge
group $SU(3) \times SU(2) \times U(1)$ is dimension-forced from
$D = 3$---color from $\mathrm{face\_pairs}(3) = 3$, weak isospin from
spinor structure, hypercharge from $U(1)$ phase.  No larger $G$ to
break, hence no topological defect, hence no GUT monopoles.

\begin{remark}
GUTs: $G \to \mathrm{SM}$ yields monopoles.  RS: SM from dimension,
no $G$, no monopoles.
\end{remark}

\section{Experimental Status of Monopole Searches}
\label{sec:experimental}

\textbf{MACRO}~\cite{MACRO} (Gran Sasso): GUT monopole search, flux
$\Phi \lesssim 10^{-16}$ cm$^{-2}$ sr$^{-1}$ s$^{-1}$.  No candidates.
\textbf{IceCube}~\cite{IceCube}: Cherenkov limits on slow/relativistic
monopoles; no signal.  \textbf{MoEDAL}~\cite{MoEDAL} (LHC): track
detectors for $pp$-produced monopoles; Run 2 null, mass $\lesssim 5$
TeV limits.

Null results across mass and velocity ranges support the RS prediction.

\section{Predictions and Falsifiers}
\label{sec:predictions}

\textbf{Predictions:} (i) No fundamental Dirac monopoles---charge
quantization is ledger-derived.  (ii) No GUT monopoles---gauge group
is dimension-forced, not embedded.  (iii) All searches (MoEDAL,
IceCube, MACRO, future) return null results.

\textbf{Falsifiers:} Discovery of any magnetic monopole would falsify
RS.  Evidence for a GUT (proton decay, coupling unification) would
challenge the claim that the SM gauge group is terminal.

\section{Conclusion}
\label{sec:conclusion}

Charge quantization in RS is dimensional and ledger-derived, not
monopole-driven.  The discrete ledger forces integer charges;
$\mathrm{face\_pairs}(3) = 3$ fixes three charge types.  Dirac's 1931
argument is circumvented.  RS predicts no monopoles---fundamental or
GUT.  The gauge group emerges from dimension alone; no larger symmetry
to break.  Experimental searches support this prediction.

\begin{thebibliography}{99}
\bibitem{Dirac1931}
P.~A.~M.~Dirac, Proc.\ R.\ Soc.\ Lond.\ A \textbf{133}, 60 (1931).

\bibitem{tHooft1974}
G.~'t Hooft, Nucl.\ Phys.\ B \textbf{79}, 276 (1974).

\bibitem{Polyakov1974}
A.~M.~Polyakov, J.\ Exp.\ Theor.\ Phys.\ Lett.\ \textbf{20}, 194 (1974).

\bibitem{MACRO}
MACRO Collaboration, Eur.\ Phys.\ J.\ C \textbf{25}, 511 (2002).

\bibitem{IceCube}
IceCube Collaboration, Eur.\ Phys.\ J.\ C \textbf{76}, 133 (2016).

\bibitem{MoEDAL}
MoEDAL Collaboration, Phys.\ Rev.\ Lett.\ \textbf{118}, 061801 (2017);
  B.~Acharya et al., Phys.\ Rev.\ Lett.\ \textbf{123}, 021802 (2019).

\bibitem{RS_Gauge}
J.~Washburn, ``The Gauge Group $SU(3) \times SU(2) \times U(1)$ from
  Recognition Science,'' RS\_Gauge\_Group\_Origin (2026).

\bibitem{RS_Axioms}
J.~Washburn, ``Recognition Science: Derivation of the Cost
Functional and Forcing Chain,'' Recognition Science Research
Institute Technical Report (2025).
\end{thebibliography}

\end{document}
